% summary.tex -- Module 7: Last-minute summary (the night before the exam).
% Facts and numbers taken from ch1.tex-ch4.tex (no new calculations).
\unnumberedchapter[\faClipboardList]{Last-minute summary}\label{bm:summary}

Read this the night before the exam. Each box gives the must-know ideas, the formulas and the
traps of one chapter. \textcolor{hscmistake}{\faTimes} marks a common mistake.

\begin{summarybox}[fontupper=\small]{Chapter 1: The electromagnetic spectrum}
\begin{itemize}[itemsep=0pt, topsep=2pt]
  \item Maxwell unified electricity and magnetism. His new idea: a changing electric field
        makes a magnetic field. He predicted EM waves moving at
        $1/\sqrt{\mu_0\varepsilon_0} = \qty{3.00e8}{m.s^{-1}}$, so light is an EM wave. Hertz
        made radio waves in 1887.
  \item EM waves come from \emph{accelerating} charges. $E$ and $B$ are in phase, at \ang{90}
        to each other and to the travel. No medium is needed; all travel at $c$ in a vacuum.
  \item Measuring $c$: Galileo (failed), Rømer, Fizeau, Foucault, Michelson, lasers. Since
        1983, $c$ is exactly \qty{299792458}{m.s^{-1}}, and it defines the metre.
  \item Spectra: hot, dense $\to$ continuous; excited thin gas $\to$ bright lines; cool gas in
        front of a hot source $\to$ dark lines at the \emph{same} wavelengths. Stars: lines
        $\to$ composition; peak $\to$ temperature; shift $\to$ motion towards or away (only
        along the line of sight); wider lines $\to$ faster rotation or higher density.
\end{itemize}
\textbf{Formulas:} $c = f\lambda$ (sheet: $v = f\lambda$), $\lambda_\mathrm{max} = b/T$
(sheet); $c = d/t$, $c = 1/\sqrt{\mu_0\varepsilon_0}$, $\Delta\lambda = \lambda v/c$ (learn).
\par\textbf{Traps:} \textcolor{hscmistake}{\faTimes} \si{nm} or \si{GHz} not converted;
\textcolor{hscmistake}{\faTimes} one-way distance for an echo ($d = 2D$);
\textcolor{hscmistake}{\faTimes} $T$ in \si{\degreeCelsius}; \textcolor{hscmistake}{\faTimes}
``Maxwell measured $c$'' (he predicted it); \textcolor{hscmistake}{\faTimes} the sheet's
Doppler formula for light (it is for sound).
\end{summarybox}

\begin{summarybox}[fontupper=\small]{Chapter 2: Light as a wave}
\begin{itemize}[itemsep=0pt, topsep=2pt]
  \item Newton: particles. Huygens: waves. Both explain reflection. Refraction: Newton needs
        light to be \emph{faster} in water, Huygens \emph{slower}. Foucault (1850): slower.
  \item Diffraction (spreading through a gap, most when the gap is near $\lambda$) and
        interference (Young, about 1801) show that light is a wave: only waves cancel.
  \item Bright fringes where $\delta = d\sin\theta = m\lambda$; dark where
        $\delta = \left(m - \frac{1}{2}\right)\lambda$. A grating gives sharp lines; red bends
        most; the highest order is $d/\lambda$ rounded \emph{down}.
  \item Only transverse waves can be polarised, so light is transverse (as Maxwell predicts).
        The first polariser passes half of unpolarised light; after that, use Malus' law.
\end{itemize}
\textbf{Formulas:} $d\sin\theta = m\lambda$, $n_\mathrm{x} = c/v_\mathrm{x}$,
$I = I_\mathrm{max}\cos^2\theta$ (sheet); $\Delta y = \lambda L/d$, $y = L\tan\theta$,
$d = 1/N$, $I = \frac{1}{2}I_0$, $\delta = |S_2P - S_1P|$ (learn).
\par\textbf{Traps:} \textcolor{hscmistake}{\faTimes} slit width used for $d$;
\textcolor{hscmistake}{\faTimes} lines per \si{mm} not $\times 1000$;
\textcolor{hscmistake}{\faTimes} $\cos\theta$ instead of $\cos^2\theta$;
\textcolor{hscmistake}{\faTimes} $\theta$ not measured from the filter before;
\textcolor{hscmistake}{\faTimes} $\Delta y = \lambda L/d$ at large angles.
\end{summarybox}

\begin{summarybox}[fontupper=\small]{Chapter 3: Light as a particle}
\begin{itemize}[itemsep=0pt, topsep=2pt]
  \item A hotter black body is brighter at every wavelength, with its peak at a shorter
        wavelength (Wien). Classical physics predicted infinite ultraviolet intensity (the
        ultraviolet catastrophe). Planck (1900): energy comes in quanta, $E = hf$.
  \item Photoelectric results: a threshold frequency; no delay; $K_\mathrm{max}$ depends on
        frequency, not brightness. The wave model cannot explain these. Brighter light gives
        more electrons, not faster ones.
  \item Einstein (1905): light is photons; one photon gives all its energy to one electron.
        Energy conservation: $K_\mathrm{max} = hf - \phi$, with $\phi = hf_0$.
  \item $K_\mathrm{max}$--$f$ graph: gradient $h$, $f$-intercept $f_0$, $K_\mathrm{max}$-intercept
        $-\phi$; $V_0$--$f$ graph: gradient $h/q$. Same gradient for every metal.
  \item Planck quantised the oscillators; Einstein quantised light. Light has wave--particle
        duality.
\end{itemize}
\textbf{Formulas:} $E = hf$, $K_\mathrm{max} = hf - \phi$, $\lambda_\mathrm{max} = b/T$
(sheet); $E = hc/\lambda$, $K_\mathrm{max} = qV_0$ (combine); $\phi = hf_0$,
$\lambda_0 = hc/\phi$, $N/t = P/(hf)$ (learn).
\par\textbf{Traps:} \textcolor{hscmistake}{\faTimes} mixing \si{eV} and \si{J}
(\qty{1}{\eV} $=$ \qty{1.602e-19}{J}); \textcolor{hscmistake}{\faTimes} a wavelength
\emph{longer} than $\lambda_0$ (it must be shorter); \textcolor{hscmistake}{\faTimes}
$K_\mathrm{max}$ in \si{eV} put into $K = \frac{1}{2}mv^2$; \textcolor{hscmistake}{\faTimes}
``the intercept is $\phi$'' (it is $-\phi$).
\end{summarybox}

\begin{summarybox}[fontupper=\small]{Chapter 4: Light and special relativity}
\begin{itemize}[itemsep=0pt, topsep=2pt]
  \item Postulate 1: the laws of physics are the same in every inertial (not accelerating)
        frame. Postulate 2: every inertial observer measures the same $c$ in a vacuum.
        Evidence: Michelson--Morley's null result, fast pions, laser tests, muons, clocks, GPS.
  \item ``At the same time'' depends on the observer. A moving clock runs slow; a moving
        object is shorter, but only \emph{along} its motion. $t_0$: the clock present at both
        events. $l_0$: measured where the object is at rest.
  \item Muons: Earth frame, lifetime dilated; muon frame, atmosphere contracted. Both agree.
        Also: Hafele--Keating, accelerators, Type Ia supernovae (stretched by $1 + z$).
  \item As $v \to c$, $p_\mathrm{v} \to \infty$, so nothing with mass reaches $c$.
  \item $E = mc^2$: the Sun loses about \qty{4.28e9}{kg} each second (fusion: about 0.7\,\% of
        the mass); annihilation gives \emph{two} photons (all the mass); burning fuel loses
        far too little mass to measure.
\end{itemize}
\textbf{Formulas:} time dilation, length contraction, $p_\mathrm{v}$ and $E = mc^2$ (all on
the sheet); $\Delta m/\Delta t = P/c^2$ (combine).
\par\textbf{Traps:} \textcolor{hscmistake}{\faTimes} $\sqrt{1 - 0.9}$ instead of
$\sqrt{1 - 0.9^2}$; \textcolor{hscmistake}{\faTimes} not checking $t > t_0$ and $l < l_0$;
\textcolor{hscmistake}{\faTimes} $v$ as a fraction of $c$ on the top line of $p_\mathrm{v}$;
\textcolor{hscmistake}{\faTimes} a mass in \si{\amu} times $c^2$ (multiply by 931.5 for
\si{\mega\eV}); \textcolor{hscmistake}{\faTimes} one photon from annihilation.
\end{summarybox}
